% from ega4-14-fr.tex
\oldpage[IV-3]{213}

\begin{corollary}[14.4.6]
\label{IV.14.4.6-fr}
Soient $Y$ un préschéma localement noethérien, $f:X\to Y$ un morphisme localement
de type fini. Soit $y$ un point géométriquement unibranche de $Y$, et supposons en
outre que pour tout $x\in f^{-1}(y)$, l'anneau $\mathcal O_{X,x}$ soit
équidimensionnel. Alors les conditions suivantes sont équivalentes~:
\begin{enumerate}[label=\emph{\alph*)}]
  \item $f$ est équidimensionnel (13.3.2) en tous les points de $f^{-1}(y)$.
  \item $f$ est ouvert en tous les points de $f^{-1}(y)$.
  \item $f$ est universellement ouvert en tous les points de $f^{-1}(y)$.
\end{enumerate}
\end{corollary}

En effet, $c)$ implique trivialement $b)$, et $a)$ implique $c)$ en vertu de
(14.4.4)~; enfin, en raison des hypothèses sur $\mathcal O_y$ et $\mathcal O_x$,
$b)$ implique $a)$ par (14.2.2). Plus généralement~:

\begin{proposition}[14.4.7]
\label{IV.14.4.7-fr}
Soient $Y$ un préschéma localement noethérien, $f:X\to Y$ un morphisme localement
de type fini, $x$ un point de $X$, tel que $y=f(x)$ soit géométriquement unibranche.
Les conditions suivantes sont équivalentes~:
\begin{enumerate}[label=\emph{\alph*)}]
  \item $f$ est équidimensionnel (13.3.2) au point $x$.
  \item L'anneau $\mathcal O_x$ est équidimensionnel et $f$ est ouvert aux points
  génériques des composantes irréductibles de $f^{-1}(y)$ contenant $x$ (donc
  aussi en tout point d'une telle composante).
  \item L'anneau $\mathcal O_x$ est équidimensionnel, et $f$ est universellement
  ouvert aux points génériques des composantes irréductibles de $f^{-1}(y)$
  contenant $x$ (donc aussi en tout point d'une telle composante).
\end{enumerate}
En outre, lorsque ces conditions sont satisfaites, pour tout sous-préschéma fermé
réduit $X_i$ de $X$ ayant pour espace sous-jacent une composante irréductible de $X$
contenant $x$, la restriction $f_i:X_i\to Y$ de $f$ est un morphisme
équidimensionnel au point $x$, et universellement ouvert en tous les points des
composantes irréductibles de $f^{-1}(y)\cap X_i$ qui contiennent $x$.
\end{proposition}

La condition $a)$ implique que $f$ est équidimensionnel en tous les points
génériques des composantes irréductibles de $f^{-1}(y)$ contenant $x$ (13.3.1), et
par suite (14.4.4) universellement ouvert en ces points~; le même raisonnement
appliqué à chaque $f_i$ (compte tenu de (13.3.3)) prouve la dernière assertion de
la proposition, compte tenu de (14.3.3.1, (iii)). En outre, d'après (14.2.1), on a
les relations
\begin{equation}
\tag{14.4.7.1}
\label{IV.14.4.7.1-fr}
  \dim(\mathcal O_{X_i,x})=\dim(\mathcal O_y)+
  \dim(\mathcal O_{X_i,x}\otimes_{\mathcal O_y}k(y))
\end{equation}
\begin{equation}
\tag{14.4.7.2}
\label{IV.14.4.7.2-fr}
  \dim(\mathcal O_{X,x})=\dim(\mathcal O_y)+
  \dim(\mathcal O_{X,x}\otimes_{\mathcal O_y}k(y))
\end{equation}
et puisque $f$ est équidimensionnel au point $x$, il résulte de (13.3.1) que l'on
a $\dim_x(X_i\cap f^{-1}(y))=\dim_x(f^{-1}(y))$, donc (5.2.3)
\[
  \dim(\mathcal O_{X_i,x}\otimes_{\mathcal O_y}k(y))=
  \dim(\mathcal O_{X,x}\otimes_{\mathcal O_y}k(y)).
\]
On conclut donc que $\dim(\mathcal O_{X_i,x})=\dim(\mathcal O_{X,x})$ pour tout
$i$, autrement dit $\mathcal O_{X,x}$ est équidimensionnel, et cela achève de
montrer que $a)$ entraîne $c)$. Il est clair que $c)$ entraîne $b)$~; enfin, $b)$
entraîne la relation (14.4.7.2) en vertu de (14.2.1)~; il résulte alors de (13.3.6)
que $b)$ entraîne $a)$.

\begin{proposition}[14.4.8]
\label{IV.14.4.8-fr}
Soient $Y$ un préschéma noethérien, $f:X\to Y$ un morphisme localement de type
fini, $y$ un point de $Y$, $x$ un point maximal de $X_y=f^{-1}(y)$. Les conditions
suivantes sont équivalentes~:
\begin{enumerate}[label=\emph{\alph*)}]
  \item Le morphisme $f$ est universellement ouvert au point $x$, autrement dit,
  pour tout changement de base $g:Y'\to Y$, on a la propriété $P(Y')$~: pour tout
  point $x'$ de $X'=X\times_Y Y'$ au-dessus de $x$, le morphisme
  $f'=f_{(Y')}:X'\to Y'$ est ouvert au point $x'$.

  